\chapter{Fichiers, commandes et captures}

\section{Fichiers principaux}

\begin{longtable}{>{\ttfamily}p{0.34\textwidth}p{0.58\textwidth}}
\caption{Inventaire fonctionnel}\label{tab:fichiers}\\
\toprule
Fichier & Fonction \\
\midrule
\endfirsthead
\toprule
Fichier & Fonction \\
\midrule
\endhead
XmiToJson.py & Conversion XMI vers JSON. \\
GrapheBuilder\_v3.py & Construction du graphe et indexation documentaire. \\
graphrag\_retriever.py & Retrieval UML, métier et documentaire. \\
Agent\_LG.py & Génération contrôlée de la réponse. \\
api.py & API, authentification, conversations et messages. \\
index.html & Interface du chatbot. \\
pipeline\_graphe\_metier.py & Extraction des propositions métier. \\
validation\_metier.py & Contrôles et corrections déterministes. \\
extraire\_tout\_graphe\_metier.py & Extraction des 29 fiches avec reprise. \\
recontroler\_graphe\_metier.py & Recontrôle hors ligne des checkpoints. \\
graphe\_metier\_builder.py & Import protégé des sélections validées. \\
importer\_pilote\_metier.py & Précontrôle et import du pilote poste. \\
evaluate\_graphrag.py & Scénarios d'évaluation fonctionnelle. \\
\bottomrule
\end{longtable}

\section{Commandes utiles}

\begin{lstlisting}[language=bash,caption={Tests et lancement de l'application}]
python -m pytest -q
python -m uvicorn api:app --host 127.0.0.1 --port 8000
\end{lstlisting}

\begin{lstlisting}[language=bash,caption={Extraction et contrôle métier}]
python extraire_tout_graphe_metier.py --inventaire
python recontroler_graphe_metier.py --output-dir DataMetier/revue_suivante
python importer_pilote_metier.py
python importer_pilote_metier.py --importer
\end{lstlisting}

\section{Requêtes Neo4j}

\begin{lstlisting}[language={},caption={Visualisation du pilote métier}]
MATCH (c:ConceptMetier {domaine: "Pilote poste d'aiguillage"})
OPTIONAL MATCH (c)-[r]->(n)
RETURN c, r, n;
\end{lstlisting}

\begin{lstlisting}[language={},caption={Visualisation de la classe Ligne}]
MATCH (c:Classe {nom: "Ligne"})
OPTIONAL MATCH (c)-[r]-(n)
RETURN c, r, n
LIMIT 40;
\end{lstlisting}

\section{Captures à ajouter}

Placer les images dans le dossier \code{figures/} avec les noms exacts indiqués dans \code{figures/README.md}. Le document compile sans elles et affiche des cadres réservés. Toute clé, tout mot de passe, cookie, jeton, identifiant personnel ou adresse interne doit être masqué avant insertion.

